\documentclass[10pt,a4paper]{amsart}
\usepackage[margin=19mm]{geometry}
\usepackage{amsmath,amssymb,mathtools}
\usepackage[scr=rsfs,cal=euler]{mathalfa}
\usepackage{xfrac}
\usepackage[unicode,hidelinks,bookmarks=false]{hyperref}
\newcommand{\ie}{i.e.\ }
\newcommand{\cf}{cf.\ }
\newcommand{\resp}{respectively\ }
\newcommand{\Subsection}{\subsection}
\providecommand{\nobreakdash}{\nobreak\hbox{-}}
\setlength{\emergencystretch}{2em}
\multlinegap0pt
\usepackage{bm}
\usepackage{bbm}
\usepackage{dsfont}
\usepackage{xspace}
\usepackage{cancel}
\usepackage{mathtools}
\usepackage{amsthm}
\makeatletter
\@ifundefined{lemma}{\newtheorem{lemma}{Lemma}}{}
\makeatother
\title[Gap distribution of $\sqrt{n} \,\mathrm{mod}\, 1$ and the ci]{Gap distribution of $\sqrt{n} \,\mathrm{mod}\, 1$ and the circle method}
\author{Maksym Radziwi{\l}{\l}, Niclas Technau}
\date{Source: arXiv:2403.16493v2 (CC BY 4.0)}
\begin{document}
\maketitle

\noindent\emph{Source and licence.} Gap distribution of $\sqrt{n} \,\mathrm{mod}\, 1$ and the circle method. Version arXiv:2403.16493v2 (CC BY 4.0), distributed under the Creative Commons Attribution 4.0 International licence, as recorded in the arXiv licence metadata for this exact version and shown on the abstract page https://arxiv.org/abs/2403.16493v2. Full rights notice: licenses/arxiv-2403.16493-CC-BY-4.0.txt.\par\smallskip

\noindent\textit{Editorial scope.} Counted unit: the Lemma whose source label is \texttt{le:paircorrel} in the article (source0.tex lines 464--479, source label \texttt{le:paircorrel}), delivered together with the article's own complete original proof of it (source0.tex lines 480--821, one \texttt{proof} environment of 342 lines). No mathematics is written, repaired or re-derived: statement and proof are the article's own text line for line. Required context retained uncounted: eq:qu (lines 254--264). Every definition, display and label of those lines is retained unchanged. Omitted from this package and not redistributed: 1--253, 265--463, 822--2024. Nothing inside a retained excerpt is omitted or altered: every retained body is delivered exactly as the article's lines are written, so no omission receipt of any kind accompanies this package. Citations: the retained excerpts carry no citation command at all, so no bibliography entry of the article is transcribed and no reference is redistributed. Reference adaptation: none was needed and none was made; every reference inside the delivered unit resolves inside it, so no unresolved reference or citation is produced. Printed slips kept literal: no printed word, symbol, hypothesis or number of the counted statement or of its proof was altered. Layout and class: the article's own class is \texttt{amsart} with packages including amsmath, amssymb, amsfonts, amscd, latexsym, amsthm, mathrsfs, verbatim, textcomp, hyperref, cite, hypbmsec, enumerate, bm; the wrapper is the lane's pinned \texttt{amsart} editorial wrapper at 10pt on A4, which restates the theorem-like environments the retained text needs and adds the article's own macro definitions that the retained text needs (none), with extra wrapper packages (bm, bbm, dsfont, xspace, cancel, mathtools). The delivered numbering is the unit's own. Assets: no retained span contains a figure environment, a graphics-inclusion command, a PDF-page inclusion, a table or a TikZ picture; no image file is copied into this package, the package declares no asset dependency and no image file is redistributed with it. Licence: version arXiv:2403.16493v2 of this article is distributed under the Creative Commons Attribution 4.0 International licence, as recorded in the arXiv licence metadata for this exact version and shown on the abstract page https://arxiv.org/abs/2403.16493v2. A full rights notice is delivered at licenses/arxiv-2403.16493-CC-BY-4.0.txt. Characters: every retained excerpt is pure ASCII, so the delivered engine, XeLaTeX with its default Latin Modern text font, reports no missing character.
    \begin{align} \label{eq:qu}
    \mathcal{Q}_{\Delta, N} := 
    \Big \{ a b \in [Q, 2Q]: Q^{1/1000} \leq a \leq 2 Q^{1/1000}, p | a b \implies p > \Delta^{2001}, \mu^2(a b) = 1 \Big \}
    % \Bigg\{ a b \in [Q, 2Q] :
    % \begin{array}{l} 
    %  Q^{1/1000} \leq a \leq 2 Q^{1/1000}  \\
    %  \ p | a b \implies p > \Delta^{2001}  \\
    %   \mu^2(a b) = 1
    %  \end{array}
    %  \Bigg\},
    \end{align}

\begin{lemma}\label{le:paircorrel}
 Let $\Delta, N \geq 10$ and set 
 $Q := \Delta \sqrt{N}$. 
  Let $\mathcal{Q}=\mathcal{Q}_{\Delta,N}$ 
  be as in \eqref{eq:qu}. Suppose $\Psi$
  is a non-negative, even, 
  smooth function that is compactly supported in 
  $[-1 / 10,1 / 10]$.
  Define $\mathcal{S}=\{a/q: 
  1\leq a < q, (a,q)=1, q\in \mathcal{Q}\}$. Then,
  \begin{equation}\label{eq:paircorrelasymp}
    R_2(\mathcal{S},\Psi,\lambda) = 
    \lambda \widehat{\Psi}(0) + O(\Delta^{-2000})  
  \end{equation}
uniformly for all $\Delta\ll \lambda \leq \Delta^{100}$ as $N \rightarrow \infty$.
\end{lemma}

\begin{proof}
To analyze $R_2(\mathcal{S},\Psi,\lambda)$,
we decompose the off-diagonal in $\mathcal{S}^2$ 
into the sets $\mathcal{S}^2(g)$ consisting of all
$(s_1,s_2)=(a_1/q_1,a_2/q_2)\in  
\mathcal{S}^2$ with $g=(q_1,q_2)$ and $s_1\neq s_2$.
Then, 
\begin{equation}\label{eq gcd decomp}
R_2(\mathcal{S},\Psi,\lambda)
= 
\frac{1}{L} 
\sum_{g\geq 1} \sum_{\substack{(\frac{a_1}{q_1},
\frac{a_2}{q_2})\in \mathcal{S}^2(g)}} 
\Psi\Big(\frac{\Vert \frac{a_1}{q_1} - 
\frac{a_2}{q_2} \Vert_{\mathbb{R}/\mathbb{Z}}}
{\lambda L^{-1}}\Big).
\end{equation}
We claim $\mathcal{S}^2(g)$ contributes zero to \eqref{eq gcd decomp} if $g>1$. To verify this,
we will show that there
\begin{equation}\label{eq lower bound on fractions}
\frac{\Vert \frac{a_1}{q_1} - 
\frac{a_2}{q_2} \Vert_{\mathbb{R}/\mathbb{Z}}}
{\lambda L^{-1}} > \Delta^{1000}
\end{equation}
is inevitably outside of 
$\mathrm{supp}(\Psi)\subseteq [-1/10,1/10]$.
Let $(a_1/q_1,a_2/q_2)\in \mathcal{S}^2(g)$. We start 
by writing $q_1=q_1' g$ and $q_2=gq_2'$ so that $(q_1',q_2')=1$.
Notice
$$
\frac{a_1}{q_1} - 
\frac{a_2}{q_2}
= 
\frac{a_1q_2' - a_2 q_1'}{g q_1'q_2'}.
$$
The left-hand side is non-vanishing (as we
are on the off-diagonal $a_1/q_1\neq a_2/q_2$).
Hence, the integer $a_1q_2' - a_2 q_1'$ is non-zero
ensuring 
$\vert a_1q_2' - a_2 q_1'\vert \geq 1$ and
$$
\Big \vert 
\frac{a_1}{q_1} - 
\frac{a_2}{q_2}
\Big \vert \geq \frac{1}{g q_1'q_2'}.
$$
By the roughness stipulation in \eqref{eq:qu}, 
if $g>1$ then $g> \Delta^{2001}$.
Note that $L\leq \Delta^5 N$ and that $q_{1}', q_{2}' \leq 2 Q / g$.
Thus, 
$$
\frac{1}{g q_1'q_2'} \geq \frac{g}{4 Q^{2}} \geq \frac{\Delta^{2001}}{4 \Delta^{2} N}
> \frac{\Delta^{1800}}{L}.
$$
As $\lambda \leq \Delta^{100}$,
we deduce \eqref{eq lower bound on fractions}.
Consequently, \eqref{eq gcd decomp} simplifies to
$$
R_2(\mathcal{S},\Psi,\lambda)
= 
\frac{1}{L}  \sum_{\substack{(\frac{a_1}{q_1},
\frac{a_2}{q_2})\in \mathcal{S}^2(1)}} 
\Psi\Bigg(\frac{\big \Vert \frac{a_1}{q_1} - 
\frac{a_2}{q_2} \big\Vert_{\mathbb{R}/\mathbb{Z}}}
{\lambda L^{-1}}\Bigg).
$$
Fix co-prime
$q_1,q_2\in \mathcal{Q}$.
Let $\mathfrak{q}:=q_1q_2$.
Given an integer $m\geq 1$, we use the abbreviation $\mathbb{Z}_{m}^*:=\{1\leq h\leq m: (h,m)=1\}$.
The map
$(a_1,a_2)\mapsto a_1 q_2 - q_1 a_2 \,\mathrm{mod}\, \mathfrak{q}$
is a bijection 
between $\mathbb{Z}_{q_1}^*\times \mathbb{Z}_{q_2}^*$
and $\mathbb{Z}_{\mathfrak{q}}^*$. Moreover if 
$ a_1 q_2 - q_1 a_2 \,\mathrm{mod}\, \mathfrak{q} = h$, then
$$
\frac{1}{\lambda L^{-1}}
\Big \Vert \frac{a_1}{q_1} - 
\frac{a_2}{q_2} \Big\Vert_{\mathbb{R}/\mathbb{Z}}
= 
\min\Big( \frac{L}{\lambda}\cdot \frac{h}{\mathfrak{q}},
 \frac{L}{\lambda}\cdot  \frac{\mathfrak{q}-h}{\mathfrak{q}} \Big).
$$
Let $\mathcal{I}:=\mathrm{supp}(\Psi) \cap \mathbb{R}_{\geq 0}$. Next, we 
%show
%$\Psi\big(\frac{\Vert h/\mathfrak{q} 
%\Vert_{\mathbb{R}/\mathbb{Z}}}{\lambda L^{-1}}\big)$
%is non-zero
%if and only if $h$ is in one of two intervals. 
%Let
%$B_-:= \inf(\mathrm{supp}(\Psi)\cap\mathbb{R}_{\geq 0})$, 
%$B_+:= \sup(\mathrm{supp}(\Psi)\cap\mathbb{R}_{\geq 0})$,
%and $\{1\leq h\leq \mathfrak{q}\}=:\mathbb{Z}_\mathfrak{q}$.
observe that
$$
\Psi \Big ( \frac{L}{\lambda} \cdot \Big \| \frac{h}{\mathfrak{q}} \Big \|_{\mathbb{R} / \mathbb{Z}} \Big ) \neq 0
$$
if and only if,
$$
\min \Big (h, \mathfrak{q} - h \Big ) \in \frac{\mathfrak{q} \lambda}{L} \cdot \mathcal{I}.
$$
Thus,
$$
\Bigg\{h \in \mathbb{Z}_\mathfrak{q}: 
\Psi\Big( \frac{L}{\lambda}
\Big \Vert \frac{h}{\mathfrak{q}} 
\Big\Vert_{\mathbb{R}/\mathbb{Z}}\Big)
\neq 0
\Bigg\}
=
\Big(\frac{\mathfrak{q} \lambda}{L}
\mathcal{I}\cup 
\big(\mathfrak{q} - 
\frac{\mathfrak{q} \lambda}{L} \mathcal{I},
\mathfrak{q}\big)\Big)
\cap \mathbb{Z}_\mathfrak{q}.
$$
Each of the two intervals on the right-hand side has length comparable to
$\lambda/5 \leq \Delta^{100}$.
Using their location, we can add co-primality conditions to
both sides of the equation:
$$
\Bigg\{h \in \mathbb{Z}_\mathfrak{q}^*: 
\Psi\Big( \frac{L}{\lambda} \Big \Vert  
\frac{h}{\mathfrak{q}} \Big\Vert_{\mathbb{R}/\mathbb{Z}}\Big)
\neq 0
\Bigg\}
=
\Big(\frac{\mathfrak{q}\lambda}{L}
\mathcal{I}\cup 
\big(\mathfrak{q} - 
\frac{\mathfrak{q} \lambda}{L}\mathcal{I},
\mathfrak{q}\big)\Big)
\cap \mathbb{Z}^*_\mathfrak{q}.
$$
To proceed, it is helpful to have more information 
on $L$ at hand. To this end, we write 
$$
L = \sum_{q \in \mathcal{Q}} \varphi(q) = \sum_{q\in \mathcal{Q}} \sum_{d\mid q} \frac{q}{d} \mu (d)
=\sum_{1 \leq  d \leq 2Q} 
 \frac{\mu (d)}{d}
\sum_{
\substack{q\in \mathcal{Q}\\
d\mid q}}  q
= 
\sum_{q\in \mathcal{Q}}  q
+ 
E
$$
where $Q = \Delta \sqrt{N}$ and the contribution $E$ (from the range $d>1$) is readily estimated via
$$
E= 
\sum_{\Delta^{2001} \leq  d \leq 2Q} 
 \frac{\mu (d)}{d}
\sum_{
\substack{q\in \mathcal{Q}\\
d\mid q}}  q\ll 
\sum_{\Delta^{2001} \leq d \leq 2Q} 
 \frac{1}{d}
\Big( \frac{Q}{d} + 1 \Big) Q
\ll \frac{Q^2}{\Delta^{2001}}.
$$
As
$\#\mathcal{Q}\asymp \frac{Q}{\log \Delta}$,
we infer 
$\sum_{q\in \mathcal{Q}}  q \asymp \frac{Q^2}{\log \Delta}
\asymp \frac{\Delta^2 N}{\log \Delta} $.
%\begin{equation}\label{eq sum q}
%\sum_{q\in \mathcal{Q}}  q \asymp \frac{Q^2}{\log \Delta}.
%\end{equation}
Consequently,
\begin{equation}\label{eq: approx to L}
L= (1+ O(\Delta^{-2000})) \sum_{q\in \mathcal{Q}}  q.
\end{equation}
By $\lambda \gg \Delta$, we infer from
\eqref{eq: approx to L} that
$\lambda \mathfrak{q}/L\gg \Delta$.
Poisson summation implies
\begin{align*}
\sum_{\substack{h\in \frac{\mathfrak{q} \lambda}{L}\mathcal{I}\\
(h,\mathfrak{q})=1}}
\Psi\Big( \frac{L}{\lambda}\Big \Vert  
\frac{h}{\mathfrak{q}} \Big\Vert_{\mathbb{R}/\mathbb{Z}}\Big)
= 
\frac{1}{2}
\sum_{h\in \mathbb{Z}}
\Psi\Big(\frac{Lh}{\lambda\mathfrak{q}} \Big)
= \frac{1}{2} 
\frac{\lambda\mathfrak{q}}{L}
\widehat{\Psi}(0)
+ O_A(\Delta^{-A})
\end{align*}
for any $A>0$.
In the same way,
\begin{align*}
\sum_{\substack{h\in (\mathfrak{q}- 
\frac{\mathfrak{q} \lambda}{L}\mathcal{I}, \mathfrak{q})\\
(h,\mathfrak{q})=1}}
\Psi\Big(\frac{L}{\lambda} \Big \Vert  
\frac{h}{\mathfrak{q}} \Big\Vert_{\mathbb{R}/\mathbb{Z}}\Big)
& = \frac{1}{2} 
\frac{\lambda\mathfrak{q}}{L}
\widehat{\Psi}(0)
+ O_A(\Delta^{-A}).
\end{align*}
The upshot is that we can approximately compute the summations over the 
$a_i$ as follows
$$
\sum_{\substack{(a_i,q_i)=1\\
i=1,2}} 
\Psi\Bigg(\frac{\big \Vert \frac{a_1}{q_1} - 
\frac{a_2}{q_2} \big\Vert_{\mathbb{R}/\mathbb{Z}}}
{\lambda L^{-1}}\Bigg) = 
\frac{\lambda \mathfrak{q}}{L} \widehat{\Psi}(0) + O_A(\Delta^{-A}).
$$
As a result, 
\begin{equation}\label{eq circle method interm}
R_2(\mathcal{S},\Psi,\lambda)
= 
\lambda \widehat{\Psi}(0) 
\sum_{\substack{(q_1,q_2)\in \mathcal{Q}^2\\
(q_1,q_2)=1}}
\frac{q_1q_2}{L^2} + O_A(\Delta^{-A}).
\end{equation}
We write
\begin{align}\label{eq interm q sum}
\sum_{\substack{(q_1,q_2)\in \mathcal{Q}^2\\
(q_1,q_2)=1}} \frac{q_1q_2}{L^2}
= 
\Bigg(
\sum_{\substack{q\in \mathcal{Q}}} \frac{q}{L}
\Bigg)^2
+ O\Bigg( 
\sum_{\Delta^{2001}< g < 2Q}
\sum_{\substack{(q_1,q_2)\in \mathcal{Q}^2\\
(q_1,q_2)=g}} \frac{q_1q_2}{L^2}
\Bigg).
\end{align}
For each $g<2Q$, there are 
$O(\frac{Q^2}{g^2})$
many $(q_1,q_2)\in \mathcal{Q}^2$ 
with $(q_1,q_2)=g$. Thus, 
$$
\sum_{\Delta^{2001}< g < 2Q}
\sum_{\substack{(q_1,q_2)\in \mathcal{Q}^2\\
(q_1,q_2)=g}} \frac{q_1q_2}{L^2}
\ll 
\frac{1}{\Delta^{2000}}.
$$
%%
%%%%% Comment begin
\iffalse
Notice the following. Suppose
$x_1\leq \ldots \leq  x_I$
are points in a compact interval $[m,M]$.
If we put $m:=x_0$ and $x_{I+1}:=M$, then
$$
\frac{1}{I}
\sum_{i\leq I}
x_n -\int_{m}^{M}
x\,dx =
\sum_{0\leq i\leq I}
\int_{x_i}^{x_{i+1}}
t - \frac{x_i}{I} \,dt.
$$
Abbreviate by 
$
\mathrm{disp}(x_i: i\leq I):=
\max_{i\leq I} (x_{i+1} - x_i)
$
the dispersion of $\{x_i: i\leq I\}$.
Thus, 
$$
\Big \vert 
\frac{1}{I}
\sum_{i\leq I}
x_n -\int_{m}^{M}
x \,dx \Big \vert \leq  
(M-m) \cdot 
\mathrm{disp}(x_i: i\leq I).
$$
Taking $x_q:=q/\sqrt{N}\in [\Delta,2\Delta]$ implies 
$$
\sum_{\substack{q\in \mathcal{Q}}} \frac{q}{\sqrt{N}} 
= 
\#\mathcal{Q} 
\cdot \Big(
\frac{3}{2} \Delta^2 + 
O(\Delta \mathrm{disp}(x_q: q\in \mathcal{Q}))
\Big).
$$
We claim that there exists $\delta >0$
so that
$\mathrm{disp}(x_q: q\in \mathcal{Q})\leq 
N^{-\delta}$
as soon as $\log N\geq \Delta^{2001}$.
It suffices to show
any shortish interval $[y,y+Q^{1-\delta}]\subset[Q,2Q]$
contains at least one $q\in \mathcal{Q}$.
To verify this, we first consider the set
$$
\mathcal{G}:=
    \{ b \in \mathbb{N}: 
    p | b \implies p > \Delta^{2001}, \mu^2(b) = 1 \}.
$$
Basic sieve theory implies
that
there exists $\delta >0$ and a constant $c(\Delta)>0$ so that
$$
\#(\mathcal{G}\cap [1,z]) = c(\Delta) z + O(z^{1-2\delta})
$$
uniformly for any $z \geq e^{\Delta}$.
Hence 
\begin{equation}\label{eq short interv}
\#(\mathcal{G}\cap [z,z+z^{1-\delta}])>0    
\end{equation}
uniformly for any $z \geq e^{\Delta}$.
Now, we find $q=ab\in \mathcal{Q}$ 
in a prescribed interval $[y,y+Q^{1-\delta}]$
by noticing that (given $a\in Q^{1/1000}$)
there is $b\in [y/a,y/a+(y/a)^{1-\delta}]\cap \mathcal{G}$
due to \eqref{eq short interv}.
\fi
%%
%%%%% Comment end
By \eqref{eq: approx to L} we know
$\sum_{\substack{q\in \mathcal{Q}}} \frac{q}{L}
= 1 + O(\Delta^{-2000})$.
Using this information in \eqref{eq interm q sum} 
yields
$$
\sum_{\substack{(q_1,q_2)\in \mathcal{Q}^2\\
(q_1,q_2)=1}} \frac{q_1q_2}{L^2}
= 
1
+ O(\Delta^{-2000}).
$$
In view of \eqref{eq circle method interm}, we conclude 
$R_2(\mathcal{S},\Psi,\lambda)= 
\lambda \widehat{\Psi}(0) + O(\Delta^{-2000})$.
\end{proof}

\end{document}
